%! Author = sriadityavedantam
%! Date = 3/30/26

\section{Group Homomorphisms and Isomorphisms}

\begin{definition}[Group Homomorphism]
    If we let $f:G\to H$ be a function between groups, then $f$ is a homomorphism if for all $a,b\in G$,
    \[
        f(ab)=f(a)f(b).
    \]
\end{definition}

\begin{definition}[Group Isomorphism]
    A map $f:G\to H$ is an isomorphism if it is a bijective homomorphism.
\end{definition}

\begin{example}
    Let
    \[
        G\coloneqq \left\{\exp\left(\frac{2k\pi i}{n}\right):k=0,\dots,n-1\right\}
    \]
    under multiplication.
    Let
    \[
        H\coloneqq (\mathbb{Z}_n,+).
    \]
    Define
    \[
        f:G\to H
    \]
    by
    \[
        f\left(\exp\left(\frac{2k\pi i}{n}\right)\right)=k\in\mathbb{Z}_n.
    \]
    Prove that $G$ and $H$ are isomorphic.
\end{example}

\begin{proposition}{
    The map
    \[
        f\left(\exp\left(\frac{2k\pi i}{n}\right)\right)=k
    \]
    defines an isomorphism
    \[
        G\cong \mathbb{Z}_n.
    \]
}
    First, $f$ is a homomorphism.
    If
    \[
        \exp\left(\frac{2k\pi i}{n}\right),\exp\left(\frac{2\ell\pi i}{n}\right)\in G,
    \]
    then
    \begin{align*}
        f\left(\exp\left(\frac{2k\pi i}{n}\right)\exp\left(\frac{2\ell\pi i}{n}\right)\right)
        &=
        f\left(\exp\left(\frac{2(k+\ell)\pi i}{n}\right)\right)\\
        &=
        k+\ell\\
        &=
        f\left(\exp\left(\frac{2k\pi i}{n}\right)\right)+
        f\left(\exp\left(\frac{2\ell\pi i}{n}\right)\right).
    \end{align*}
    So $f$ is a homomorphism.
    To check injectivity, note that
    \[
        \exp\left(\frac{2k\pi i}{n}\right)=\exp\left(\frac{2\ell\pi i}{n}\right)
    \]
    if and only if
    \[
        \exp\left(\frac{2(k-\ell)\pi i}{n}\right)=1,
    \]
    if and only if
    \[
        \frac{2(k-\ell)\pi}{n}
    \]
    is a multiple of $2\pi$,
    if and only if
    \[
        n\mid (k-\ell),
    \]
    if and only if
    \[
        k\equiv \ell \mod n.
    \]
    Thus $f$ is well-defined and injective.
    It is surjective because every class $k\in\mathbb{Z}_n$ is hit by
    \[
        \exp\left(\frac{2k\pi i}{n}\right).
    \]
    Hence $f$ is an isomorphism.
\end{proposition}

\begin{example}
    The map
    \[
        f:(\mathbb{Z},+)\to (\mathbb{Z}_n,+)
    \]
    defined by
    \[
        f(k)=[k]
    \]
    is a homomorphism because
    \[
        [k+\ell]=[k]+[\ell].
    \]
    It is not injective because
    \[
        f(0)=f(n)=f(2n).
    \]
    But it is surjective, since every class in $\mathbb{Z}_n$ is of the form $[k]$.
\end{example}

\begin{example}
    Let
    \[
        G\coloneqq \left\{\exp\left(\frac{2\pi ki}{n}\right):k\in\mathbb{Z}\right\}.
    \]
    Then
    \[
        f:(\mathbb{Z},+)\to G
    \]
    defined by
    \[
        f(k)=\exp\left(\frac{2\pi ki}{n}\right)
    \]
    is a homomorphism since
    \[
        f(k+\ell)=f(k)f(\ell).
    \]
\end{example}

\begin{proposition}[Isomorphism of Infinite Cyclic Groups]{
    Suppose $G$ is an infinite cyclic group.
    Then
    \[
        G\cong (\mathbb{Z},+).
    \]
}
    Given that $G$ is an infinite cyclic group under multiplication, there exists $a\in G$ such that
    \[
        G=\{a^n:n\in\mathbb{Z}\}.
    \]
    Define
    \[
        f:\mathbb{Z}\to G
    \]
    by
    \[
        f(n)=a^n.
    \]
    Then
    \[
        f(n+m)=a^{n+m}=a^{n}a^m=f(n)f(m),
    \]
    so $f$ is a homomorphism.
    If $f(m)=f(n)$, then
    \[
        a^m=a^n,
    \]
    so
    \[
        a^{m-n}=1.
    \]
    Since $G$ is infinite cyclic, $a$ has infinite order.
    Thus
    \[
        m-n=0,
    \]
    so
    \[
        m=n.
    \]
    Hence $f$ is injective.
    It is surjective because $G$ is generated by $a$.
    Therefore $f$ is an isomorphism.
\end{proposition}

\begin{proposition}[Isomorphism of Finite Cyclic Groups]{
    Suppose $G$ is a finite cyclic group under multiplication and $\ord(G)=n$.
    Then
    \[
        G\cong (\mathbb{Z}_n,+).
    \]
}
    Given $\ord(G)=n$ and $G$ is generated by $c$, define
    \[
        f:(\zn,+)\to G
    \]
    by
    \[
        f([a])=c^a.
    \]
    This is well-defined because if
    \[
        a\equiv b\mod n,
    \]
    then
    \[
        n\mid (a-b),
    \]
    so
    \[
        c^{a-b}=1,
    \]
    hence
    \[
        c^a=c^b.
    \]
    Now
    \begin{align*}
        f([a]+[b]) &= f([a+b])\\
        &= c^{a+b}\\
        &= c^{a}c^b\\
        &= f([a])f([b]).
    \end{align*}
    Thus $f$ is a homomorphism.
    If $f([a])=f([b])$, then
    \[
        c^a=c^b,
    \]
    so
    \[
        c^{a-b}=1.
    \]
    Since $\ord(c)=n$, this implies
    \[
        n\mid (a-b),
    \]
    hence
    \[
        [a]=[b].
    \]
    So $f$ is injective.
    It is surjective since every element of $G$ is of the form $c^a$.
    Therefore
    \[
        (\zn,+)\cong G.
    \]
\end{proposition}

\begin{example}
    If $p$ is prime, then
    \[
        (\mathbb{Z}_p^*,\cdot)\cong (\mathbb{Z}_{p-1},+).
    \]
\end{example}

\begin{proposition}{
    If $p$ is prime, then
    \[
        (\mathbb{Z}_p^*,\cdot)\cong (\mathbb{Z}_{p-1},+).
    \]
}
    Since $\mathbb{Z}_p^*$ is cyclic of order $p-1$, the previous proposition applies.
    Hence
    \[
        \mathbb{Z}_p^*\cong \mathbb{Z}_{p-1}.
    \]
\end{proposition}

\begin{definition}[Automorphism]
    An isomorphism
    \[
        f:G\to G
    \]
    is called an automorphism.
\end{definition}

\begin{theorem}[Properties of Homomorphisms]{
    Let $f:G\to H$ be a group homomorphism.
    Then:
    \begin{enumerate}
        \item $f(1_G)=1_H$.
        \item If $a\in G$, then $f(a^{-1})=f(a)^{-1}$.
        \item
        \[
            \ima(f)\coloneqq \{h\in H:\exists g\in G,\ f(g)=h\}
        \]
        is a subgroup of $H$.
        \item If $f$ is injective, then
        \[
            \overline{f}:G\to \ima(f)
        \]
        defined by
        \[
            \overline{f}(g)=f(g)
        \]
        is an isomorphism.
    \end{enumerate}
}
    \textbf{1.}
    Since $f$ is a homomorphism,
    \[
        f(1_G)=f(1_{G}1_G)=f(1_G)f(1_G).
    \]
    Multiply on the left by $f(1_G)^{-1}$ in $H$ to get
    \[
        1_H=f(1_G).
    \]

    \textbf{2.}
    Since
    \[
        aa^{-1}=1_G,
    \]
    we have
    \[
        1_H=f(1_G)=f(aa^{-1})=f(a)f(a^{-1}).
    \]
    Thus
    \[
        f(a^{-1})=f(a)^{-1}.
    \]

    \textbf{3.}
    Let $h_1,h_2\in \ima(f)$.
    Then there exist $g_1,g_2\in G$ such that
    \[
        f(g_1)=h_1
        \quad\text{and}\quad
        f(g_2)=h_2.
    \]
    Then
    \[
        h_{1}h_2=f(g_1)f(g_2)=f(g_{1}g_2)\in \ima(f).
    \]
    Also, if $h\in \ima(f)$, say $h=f(g)$, then by part (2),
    \[
        h^{-1}=f(g)^{-1}=f(g^{-1})\in \ima(f).
    \]
    So $\ima(f)\leq H$.

    \textbf{4.}
    If $f$ is injective, then $\overline{f}:G\to \ima(f)$ is surjective by definition of image.
    It is injective because if
    \[
        \overline{f}(g_1)=\overline{f}(g_2),
    \]
    then
    \[
        f(g_1)=f(g_2),
    \]
    and injectivity of $f$ gives
    \[
        g_1=g_2.
    \]
    Since $\overline{f}$ is clearly a homomorphism, it is an isomorphism.
\end{theorem}

\begin{theorem}[Cayley's Theorem]{
    Every group $G$ is isomorphic to a subgroup of a symmetric group, in fact to a subgroup of the group of permutations of $G$ as a set.
}
    Let $G$ be a group.
    For each $g\in G$, define
    \[
        f_g:G\to G
    \]
    by
    \[
        f_g(x)=gx.
    \]
    Each $f_g$ is a permutation of $G$, since its inverse is $f_{g^{-1}}$.
    Define
    \[
        f:G\to \sym(G)
    \]
    by
    \[
        f(g)=f_g.
    \]
    Then $f$ is a homomorphism because
    \[
        f_{gh}(x)=(gh)x=g(hx)=f_g(f_h(x)).
    \]
    So
    \[
        f(gh)=f(g)\circ f(h).
    \]
    We claim that $f$ is injective.
    Suppose
    \[
        f_g=f_h.
    \]
    Then for all $x\in G$,
    \[
        gx=hx.
    \]
    In particular, taking $x=1$ gives
    \[
        g=h.
    \]
    Thus $f$ is injective.
    By the earlier theorem,
    \[
        G\cong \ima(f),
    \]
    and $\ima(f)\leq \sym(G)$.
\end{theorem}

If $G$ is finite and $\ord(G)=n$, then
\[
    \sym(G)\cong S_n.
\]
So $G$ is isomorphic to a subgroup of $S_n$.
Note that
\[
    \ord(S_n)=n!.
\]
This is called the Cayley representation.
We learned the two-line notation of permutations.
Now let us learn cycle notation.
If
\[
    \sigma=
    \begin{pmatrix}
        1&2&3&4&5&6\\
        2&5&1&3&6&4
    \end{pmatrix},
\]
then its cycle notation is
\[
    (1\ 2\ 5\ 6\ 4\ 3).
\]
If
\[
    \begin{pmatrix}
        1&2&3&4&5&6\\
        3&4&1&6&5&2
    \end{pmatrix},
\]
then its cycle notation is
\[
    (1\ 3)(2\ 4\ 6),
\]
and we usually omit fixed points such as $(5)$.

\begin{example}
    If
    \[
        \sigma=(132)(465),
        \qquad
        \tau=(23)(56),
    \]
    then
    \[
        \sigma\tau=(13)(46).
    \]
\end{example}

\begin{example}
    If
    \[
        \sigma=(1345)\in S_5,
    \]
    then
    \[
        \sigma^2=(14)(35).
    \]
\end{example}

\begin{definition}[Conjugate]
    Let $g,x\in G$.
    Then
    \[
        gxg^{-1}
    \]
    is called the conjugate of $x$ by $g$.
\end{definition}

\begin{example}
    In $S_3$, let
    \[
        \sigma=(123),
        \qquad
        \tau=(12).
    \]
    Then
    \[
        \tau\sigma\tau^{-1}=\tau\sigma\tau=(132).
    \]
\end{example}

\begin{example}
    If
    \[
        \sigma=(13546),
        \qquad
        \tau=(245),
    \]
    then
    \[
        \tau^{-1}=(254),
    \]
    and
    \[
        \tau\sigma\tau^{-1}=(13256).
    \]
\end{example}

Conjugation preserves cycle structure.
In fact,
\[
    \tau(a_1\ a_2\ \dots\ a_k)\tau^{-1}
    =
    (\tau(a_1)\ \tau(a_2)\ \dots\ \tau(a_k)).
\]

\begin{definition}[Faithful]
    If
    \[
        g\mapsto f_g
    \]
    is a homomorphism and is injective, we call the representation faithful.
\end{definition}

\begin{definition}[Group of Automorphisms of $G$]
    \[
        \aut(G)\coloneqq \{\phi:G\to G:\phi\text{ is an automorphism of }G\}
    \]
    under composition.
\end{definition}

\begin{example}
    Define
    \[
        \phi:G\to \aut(G)
    \]
    by
    \[
        g\mapsto \phi_g,
        \qquad
        \phi_g(x)=gxg^{-1}
    \]
    for $x\in G$.
\end{example}

\begin{proposition}{
    For each $g\in G$, the map
    \[
        \phi_g(x)=gxg^{-1}
    \]
    is an automorphism of $G$, and the map
    \[
        \phi:G\to \aut(G),
        \qquad
        g\mapsto \phi_g
    \]
    is a homomorphism.
}
    We claim $\phi_g\in \aut(G)$.
    For $x,y\in G$,
    \begin{align*}
        \phi_g(xy)
        &=
        gxyg^{-1}\\
        &=
        (gxg^{-1})(gyg^{-1})\\
        &=
        \phi_g(x)\phi_g(y).
    \end{align*}
    So $\phi_g$ is a homomorphism.
    It is bijective because its inverse is
    \[
        \phi_{g^{-1}}.
    \]
    Indeed,
    \[
        (\phi_g\circ \phi_{g^{-1}})(x)
        =
        \phi_g(g^{-1}xg)
        =
        g(g^{-1}xg)g^{-1}
        =
        x,
    \]
    and similarly
    \[
        \phi_{g^{-1}}\circ \phi_g=\id_G.
    \]
    Now let $g,h\in G$.
    For $x\in G$,
    \begin{align*}
        \phi_{gh}(x)
        &=
        ghx(gh)^{-1}\\
        &=
        g(hxh^{-1})g^{-1}\\
        &=
        \phi_g(\phi_h(x)).
    \end{align*}
    So
    \[
        \phi_{gh}=\phi_g\circ \phi_h.
    \]
    Hence $g\mapsto \phi_g$ is a homomorphism.
\end{proposition}

Suppose $G$ is abelian.
Then for any $x,g\in G$,
\[
    \phi_g(x)=gxg^{-1}=x.
\]
So $\phi_g$ is the identity map for all $g\in G$.

\begin{example}
    Let
    \[
        \inn(G)\coloneqq \{\phi_g:g\in G\},
    \]
    where
    \[
        \phi_g(x)=gxg^{-1}.
    \]
    This is called the group of inner automorphisms of $G$.
    Prove
    \[
        \inn(G)\leq \aut(G).
    \]
\end{example}

\begin{proposition}{
    \[
        \inn(G)\leq \aut(G).
    \]
}
    We have already shown that each $\phi_g\in \aut(G)$, so
    \[
        \inn(G)\subseteq \aut(G).
    \]
    Since
    \[
        g\mapsto \phi_g
    \]
    is a homomorphism, we have
    \[
        \inn(G)=\ima(\phi),
    \]
    which is a subgroup of $\aut(G)$.
\end{proposition}

\begin{example}
    Compute
    \[
        \inn(S_3).
    \]
\end{example}

\begin{example}
    We have
    \[
        S_3=\{e,(12),(13),(23),(123),(132)\}.
    \]
    Let
    \[
        g=(12).
    \]
    Then
    \begin{gather*}
        \phi_g(e)=e,
        \qquad
        \phi_g((12))=(12),\\
        \phi_g((13))=(12)(13)(12)=(23),
        \qquad
        \phi_g((23))=(12)(23)(12)=(13),\\
        \phi_g((123))=(12)(123)(12)=(132),
        \qquad
        \phi_g((132))=(123).\\
    \end{gather*}
    So conjugation permutes the nonidentity elements of $S_3$ in the expected way.
\end{example}

In fact in this case, we can check that
\[
    \inn(S_3)\cong \aut(S_3).
\]